% 本文件由 LaTeX 排版 Agent 根据有效 translation.json 自动生成。
% author=Tertullian; work_id=0303; title=论表演

\chapter{论表演}

% structure: p0001 source=h1 target=omitted reason=page-title-already-rendered-by-parent

% structure: p0002 source=h2 target=section reason=relative-heading-rank; base=h2

\section{第一章}

你们天主的仆人，即将走近天主，以便向祂庄严奉献自己，务要透彻理解信德的条件、真理的缘由、基督徒纪律的法规，这些法规禁止世间的其他罪过，尤其是公开表演的娱乐。你们曾宣认并承认自己已经如此行了，请回顾这一主题，以免因真正的无知或故意的无知而犯罪。因为世俗的娱乐有如此大的力量，为了保留继续享乐的机会，它设法延长甘愿的无知，并贿赂知识使其扮演不诚实的角色。也许你们中有人被外教人的观点所引诱，他们在这件事上惯常用以下论点来逼迫我们：（1）我们通过外在事物所获得的耳目精妙享受，丝毫不与心灵和良心内的宗教相悖；（2）在任何人类的享受中，只要在其自身的时空，同时向天主奉上应有的尊荣和敬畏，任何娱乐都肯定不会冒犯天主。但这正是我们准备要证明的：这些事物与真宗教和对真天主的真顺从不相容。有些人认为，基督徒，一群随时准备赴死的人，之所以被训练成他们所实践的节制，其唯一目的就是为了使轻视生命变得不那么困难，仿佛切断了与生命的联系。\Emph{他们视之为一种艺术} 熄灭一切欲望，就他们而言，他们已使之空虚，不再有任何值得渴望之物；因而这被认为更属于人类的计划和远见，而非明确由神圣命令所规定。诚然，如果基督徒在继续享受如此巨大的娱乐的同时为天主而死，那将是一件可悲的事！事实并非如他们所说；不过，若是如此，即使基督徒的顽固也大可以完全服从如此合宜的计划、如此卓越的规矩。\par

% structure: p0004 source=h2 target=section reason=relative-heading-rank; base=h2

\section{第二章}

再者，每个人都随时准备提出这样的论点：如我们所教导，万有皆由天主创造，并赐予人供其使用，且它们必定是善的，因为都出自如此良善的根源；但在这些事物中，发现了公开表演的各种组成元素，诸如马匹、狮子、体力、乐声。那么，凡凭天主自己的创造意志而存在的事物，不能被认为是与祂相异或敌对的；若与祂不相抵触，就不能被视为对祂的敬拜者有害，正如它们确实不与祂的敬拜者相异一般。毫无疑问，甚至与公共娱乐场所相连的建筑物本身，由岩石、石头、大理石、柱子构成，也都是天主的事物，祂将各种这些东西赐予大地作为点缀；不仅如此，连场景本身也在天主自己的天空下上演。人的智慧在她自己看来是多么巧妙的辩护者，尤其是当她害怕失去任何快乐——任何世俗人生甜蜜的享受时！事实上，你会发现不少人，使他们远离我们的并非生命危险，而是享乐的危险。因为即使软弱的人对死亡也没有强烈的恐惧，\Emph{作为一笔债}他自知是欠了的；而智者并不轻视快乐，视之为宝贵的恩赐——事实上，是人生唯一的福乐，无论对哲人还是愚人。没有人否认无人不知的事——因为自然本身就是这事的教师——即天主是宇宙的造物主，宇宙是善的，并由造物主作为恩赐白白赐予人。但由于与至高者没有亲密的认识，仅凭自然启示认识祂，而非作为祂的\Quote{朋友们}——远远地，而非作为那些被带到祂面前的人——人必然对祂在治理祂的世界时所命令和所禁止的一无所知。他们也必定不知道那敌对的力量，这力量对抗祂，并扭曲祂手所造之物用于错误的目的；因为你无法认识你所不认识的天主的旨意或仇敌。那么，我们不可仅考虑万有由谁造成，还要考虑它们被谁败坏了。当我们发现它们最初为何而造时，我们将找出它们未被用于什么目的。败坏状态与原初纯洁之间有天壤之别，正是因为造物主与败坏者之间有天壤之别。为什么？各式各样的邪恶，即使是外教人也无疑视为邪恶而禁止并防范的，都出自天主的工作。例如，杀人，无论用铁器、毒药还是魔法咒语。铁、草药和魔鬼都同样是天主的受造物。难道造物主竟为人的毁灭提供了这些东西吗？不，祂以那一条总诫命禁止一切杀人行为：\BibleQuote{你不可杀人。}此外，除了天主——世界的造物主——谁将金、铜、银、象牙、木材以及制造偶像所用的所有其他材料置于其中呢？然而祂这样做，是为了让人设立一种反对祂自己的敬拜吗？相反，偶像崇拜在祂眼中是最大的罪。有什么得罪天主的事不是出于天主呢？但一旦得罪祂，它就不再是祂的；而不再属于祂时，在祂眼中就成了得罪祂的事。人本身，尽管犯有各种不义，不仅是天主的工作——他是祂的肖像，然而在灵魂和肉体上，他都已与他的造物主分离。因为我们生眼睛不是为了满足情欲，舌头不是为了口出恶言，耳朵不是为了接收恶言，喉咙不是为了服事饕餮之恶，肚腹不是为了作饕餮的同盟，生殖器不是为了不洁的放纵，手不是为了暴力行为，脚不是为了走上迷途；或者灵魂置于身体之内，是为了成为阴谋、欺诈和不义的思想工厂？我想不是；因为如果天主，作为公义的无辜要求者，憎恶一切类似恶毒的事物——如果祂完全憎恶这样的邪恶谋划，那么毫无疑问，在祂手所造的一切中，祂没有造任何一件导向祂所谴责的行为，即使这些行为可能由祂所造之物实施；因为事实上，被谴责的唯一根据，就是受造物误用了造化。因此，我们这些在对主的认识中，也获得了对祂仇敌的一些认识——在发现造物主的同时，也抓住了那大败坏者——既不应惊奇也不应怀疑：正如那败坏并敌视天主的天使起初推翻了人的美德——人本是天主的工作和肖像、世界的拥有者——同样，他也完全改变了人的本性——人的本性原本像他自己的本性一样被造为完全无罪——变成了他自己邪恶敌对造物主的状态，以致在那件因其赐予人而未赐予他而使他悲痛的事物上，他可能使人在天主眼中成为有罪的，并建立自己的至上统治。\par

% structure: p0006 source=h2 target=section reason=relative-heading-rank; base=h2

\section{第三章}

凭此知识抵御了外邦人的观点，我们不妨转而思考我们同胞中那些无价值的议论；因为有些人的信德，或因过于单纯，或因过于谨慎，要求圣经给予明确权威才肯放弃表演，并声称此事尚有疑问，因为这种戒绝并未清楚而明确地加于天主的仆人身上。然而，我们从未发现它以同样的精确性表达为：\Quote{你不可进入竞技场或剧场，你不可观看格斗或表演；}正如它清楚规定的：\Quote{你不可杀人；你不可崇拜偶像；你不可犯奸淫或欺诈。}\BibleRef{出 20:14}但是，我们发现\Person{达味}的这第一句话正与此事相关：\Quote{有福的，}他说，\Quote{是那未进入不敬虔者的集会，未站立在罪人的道路，未坐在讥讽者的座位的人。}虽然他似乎在先前预言了那位义人，说他未参与犹太人的集会与商议，即关于杀害我主的会议，然而神圣的圣经总有深远的应用：在直接意义用尽之后，它从各个方面坚固敬虔生活的实践，因此在这里你也有一个语句，与明确禁止表演相差不远。如果他称那几个犹太人为不敬虔者的集会，那么他将如何更厉害地称如此众多的外邦人集会呢？难道外邦人比当时的犹太人不敬虔、更少罪、更少基督的仇敌吗？而且，看其他事情也相符。因为在表演中，他们也站立在道路中。他们称半圆形剧场中环绕座位的间隔以及分隔向下人群的通道为\Quote{道路}。曲线处主妇们坐的位置称为\Quote{座位}。因此，相反地，这成立：不福的是那进入任何恶人的议会，站立在任何罪人的道路，坐在任何讥讽者的座位上的人。我们可以理解一件事是被普遍说出的，即使它需要某种特殊解释。因为有些话带有特殊指向，却包含普遍真理。当天主告诫以色列人他们的本分，或严厉责备他们时，祂一定指向所有人；当祂威胁毁灭\Place{埃及}和\Place{埃塞俄比亚}时，祂肯定预先谴责每一个犯罪的民族。如果，从\Emph{种}到\Emph{属}推理，每一个犯罪反对他们的民族都是\Place{埃及}和\Place{埃塞俄比亚}；那么同样，从属到种推理，关于表演的起源，每一个表演都是恶人的集会。\par

% structure: p0008 source=h2 target=section reason=relative-heading-rank; base=h2

\section{第四章}

恐怕有人认为我们仅在玩弄辩论技巧，我将转而求助于我们\Quote{印记}本身的最高权威。当我们进入水中时，我们按照其规条的话语宣认基督信仰；我们公开见证我们已经弃绝了魔鬼、他的炫耀及其使者。难道魔鬼及其炫耀和使者不是首先与偶像崇拜相关吗？简言之——因为我不想详述——由此你得到一切不洁与邪恶的灵。因此，如果证明表演的全部装置都基于偶像崇拜，那么毫无疑问，这将导致这样的结论：我们在圣洗之池中的弃绝见证是与表演有关的，这些表演通过其偶像崇拜已被交付给魔鬼、他的炫耀及其使者。于是我们将陈述它们的几个起源，在何种养育之所它们长大成人；接着是其中一些的名称，它们被称为什么；然后是它们的装置，以何种迷信被奉行；（然后是它们的场所，被奉献给何种守护者；）然后是服务于它们的技艺，被追溯到何种作者。如果发现其中任何一项与偶像神祇无关，那么它将立即被视为免于偶像崇拜的玷污，并且不属于我们圣洗弃绝的范围。\par

% structure: p0010 source=h2 target=section reason=relative-heading-rank; base=h2

\section{第五章}

关于它们的起源，由于这些事对许多人而言颇为隐晦且鲜为人知，我们的探究必须追溯到远古时代，而我们的权威典籍只能是外教文献。现有若干作者存世，他们就这个主题出版了著作。他们给出的竞技起源如下。\Person{蒂迈欧}告诉我们，来自亚细亚的移民，在\Person{蒂尔赫努斯}的率领下，定居在\Place{埃特鲁里亚}；蒂尔赫努斯曾因其本土王国之争而败于其兄弟。此外，在宗教名义下的其他迷信习俗中，他们在新家园设立了公开表演。罗马人应自己的请求，从他们那里获得了熟练的表演者、合适的节期，以及名称；据说它们被称为\OriginalTerm{竞技}{Ludi}，源自\OriginalTerm{吕底亚人}{Lydi}。尽管\Person{瓦罗}将\OriginalTerm{竞技}{Ludi}的名称源自\OriginalTerm{游戏}{Ludus}，即来自玩耍，正如他们也称\OriginalTerm{卢佩尔卡里节}{Luperci}为\OriginalTerm{卢迪伊}{Ludii}，因为他们奔跑嬉戏；然而在他看来，年轻人的那种嬉戏属于节日、庙宇及宗教敬拜的对象。然而，名字的起源无关紧要，既然可以肯定这东西源自偶像崇拜。\OriginalTerm{解放节}{Liberalia}，在\OriginalTerm{竞技}{Ludi}的通称下，清楚宣告了巴克斯父神的荣耀；因为最初是感激的农民将这些庆典奉献给巴克斯，作为他赐予他们恩惠的回报，依他们所说，揭示了葡萄酒的愉悦。然后\OriginalTerm{康苏斯节}{Consualia}被称为\OriginalTerm{竞技}{Ludi}，起初是献给尼普顿的，因为尼普顿也名为康苏斯。此后，\Person{罗慕路斯}将\OriginalTerm{骑马赛}{Equiria}献给玛尔斯，尽管他们也声称康苏斯节属于罗慕路斯，理由是他将之祝圣给康苏斯——如他们所愿，是计谋之神；当然，正是那个计谋，他策划掳掠萨宾少女作为士兵的妻子。这计谋确实高明；而且我想，罗马人自己仍将其视为公正与正直，我不敢说天主也如此认为。这也玷污了起源：你当然不能认为源自罪、无耻、暴力、仇恨、谋杀兄弟的创立者、玛尔斯之子的事物是好的。即使在今天，竞技场的第一转弯处，有同一康苏斯的地下祭台，带有铭文如下：\Quote{康苏斯，计谋伟大，玛尔斯，战斗中的强大守护神祇。}国家司铎在七月望日在此祭献；罗慕路斯的司铎与灶贞女在九月朔日前第十二日。除此之外，根据\Person{皮索}传给我们的记载，罗慕路斯在\Place{塔尔佩亚山}设立了纪念朱庇特·费雷特里乌斯的竞技，既称为塔尔佩亚竞技，也称为卡比托利竞技。在他之后，\Person{努马·蓬皮利乌斯}设立了献给玛尔斯和罗比戈（因为他们还发明了一位锈蚀女神）的竞技；然后是\Person{图卢斯·霍斯提利乌斯}；然后是\Person{安库斯·马尔提乌斯}；以及其他陆续效仿的诸多人等。至于这些竞技是为纪念哪些偶像而设立的，在\Person{苏埃托尼乌斯·特兰奎卢斯}的著作中有充分记载。但我们无需再多说以证明偶像崇拜起源的指控。\par

% structure: p0012 source=h2 target=section reason=relative-heading-rank; base=h2

\section{第六章}

古代的证据之外，再加上后来相继设立的竞技，它们从其至今仍持有的名称中暴露了起源，仿佛铭刻在其脸上，表明竞技（无论何种）是为哪个偶像、哪个宗教对象而设计的。你有以\OriginalTerm{大神母}{the great Mother}、\OriginalTerm{阿波罗}{Apollo}、\OriginalTerm{刻瑞斯}{Ceres}、\OriginalTerm{尼普顿}{Neptune}、\OriginalTerm{朱庇特·拉提亚里斯}{Jupiter Latiaris}和\OriginalTerm{福罗拉}{Flora}命名的节日，全为一个共同目的庆祝；其他的宗教起源在于国王的诞辰与庆典、公共胜利以及市政节日。还有遗嘱指定的表演，其中向私人追念致以葬礼荣誉；这是依照古代的制度。因为从一开始，\Quote{竞技}就被视为两种：神圣的和葬礼的，即纪念外邦神祇和死者。但在偶像崇拜一事上，对我们而言，无论以什么名称或头衔行之，都没有区别，既然它涉及我们所弃绝的邪恶神祇。如果向死者表达敬意是合法的，那么向他们的神祇表达敬意也同样合法：两者起源相同；偶像崇拜相同；我们方面对所有偶像崇拜的庄严弃绝也相同。\par

% structure: p0014 source=h2 target=section reason=relative-heading-rank; base=h2

\section{第七章}

因此，两种公共竞技有一个起源，并且拥有共同的名称，如同具有相同的出身。同样，既然它们同样被偶像崇拜之罪（它们的母体）所玷污，它们的排场也必然相似。但竞技游戏中更为盛大的序幕展示（即游行专指的部分）本身即证明整个活动属于谁：众多偶像、长长一排雕像、各类战车、宝座、冠冕、服饰。此外还有何等崇高的宗教礼仪，何等的祭献在先、居中及随行。多少社团、多少司铎职、多少职务被搅动，居住在那魔鬼集会总部所在的大城中的居民是知道的。如果在行省中这些事以较简朴的方式举行，依据其较低下的手段，但所有竞技游戏仍必须被视为属于它们所源出之物；它们所涌出的泉源玷污了它们。涓涓细流从源头起，嫩枝从发芽起，就含有其起源的本质。它可以宏大或卑微，无关紧要，任何竞技游行都冒犯天主。即使只有少数偶像装饰它，其中就有偶像崇拜；即使只有一辆神车，那也是朱庇特的战车：任何偶像崇拜之物，无论装备寒酸或适度富丽堂皇，都从其起源上玷污了它。\par

% structure: p0016 source=h2 target=section reason=relative-heading-rank; base=h2

\section{第八章}

依我计划中关于场所的部分：竞技场主要奉献给太阳（the Sun），太阳的神殿位于竞技场中央，其神像从神殿顶端闪耀光芒；因为他们认为，对一个暴露在露天中的对象，不宜在屋顶下献上神圣的敬意。那些断言第一个表演是由\Person{喀耳刻}为了尊崇其父太阳而举办的人，也主张竞技场之名源自她。显然，这位女巫是以她作为女祭司的神祇——我指的是魔鬼和邪灵——之名行事的。因此，你看，这场所的装饰聚集了何等多的偶像崇拜！竞技场的每一件装饰本身就是一座神殿。那些蛋被视为奉献给\OriginalTerm{卡斯托尔兄弟}{Castors}，这些人毫不羞耻地宣称它们产自一只天鹅的蛋，而那只天鹅不是别人，正是朱庇特（Jupiter）本人。海豚为尊崇\OriginalTerm{尼普顿}{Neptune}而喷水。所谓\OriginalTerm{塞西亚}{Sessia}（即播种女神）、\OriginalTerm{墨西亚}{Messia}（即收获女神）、\OriginalTerm{图图利娜}{Tutulina}（即护果之神）的雕像布满了柱子。在这些雕像前方，有三座祭坛分别献给三位神祇——\OriginalTerm{伟大、全能、胜利}{Great, Mighty, Victorious}。他们视这些神祇来自萨莫色雷斯。据\Person{赫梅特莱斯}所述，巨大的方尖碑被公开竖立献给太阳；其铭文和源头一样，属于埃及的迷信。没有他们的\OriginalTerm{伟大母亲}{Mater Magna}，魔鬼集会就黯然失色；因此她在那里的尤里普斯河上主事。\OriginalTerm{孔苏斯}{Consus}，如我们已提及的，隐藏在穆尔基亚终点柱下的地底。这两者都源自一尊偶像。因为他们认为\OriginalTerm{穆尔基亚}{Murcia}是爱神；在那里，他们为她建了一座神殿。看啊，基督徒，多少不洁之名占据了竞技场！你与一个被如此众多魔鬼充满的圣地毫无关系。说到场所，此时正好预先回答一些人的质疑。那么，你说：当游戏不举行时，我若去竞技场，会有被玷污的危险吗？我们并无法律禁止单纯进入这些场所。因为天主的仆人不仅可以进入表演场所，甚至可以进入神殿，而对其信仰毫无危险，只要他出于正当理由，与它们本来的事务和职务无关。难道街道、市场、浴室、酒馆，甚至我们的住所，不也并非全然没有偶像吗？撒殚和他的使者充满了整个世界。然而，我们背离天主，并非仅仅因为身处现世，而是因为沾染并玷染了现世的罪恶。我将与我的造物主决裂，也就是说，若我去卡皮托利诺山或塞拉比斯神殿献祭或敬拜，同样，若我去竞技场或剧院当观众也是如此。场所本身并不污染人，而是其中所行之事；甚至这些场所本身，我们认为也因此而被玷污。被玷污的事物污染我们。正是因此，我们向你们陈明这类场所是奉献给谁的，以证明其中所行之事属于那些偶像守护者，而这些场所本身正是奉献给他们的。\par

% structure: p0018 source=h2 target=section reason=relative-heading-rank; base=h2

\section{第九章}

现在谈谈竞技场表演特有的种类。昔日马术以简单方式在马上进行，其日常用途当然毫无罪过；但一旦被引入比赛，就从服务于天主转为服务于魔鬼。因此，这种竞技场表演被视为献给\OriginalTerm{卡斯托尔和波鲁克斯}{Castor and Pollux}，据\Person{斯特西科罗斯}告诉我们，马匹是由\OriginalTerm{墨丘利}{Mercury}赐给他们的。还有\OriginalTerm{尼普顿}{Neptune}也是一位马神，被希腊人称为\OriginalTerm{马神希庇乌斯}{Hippius}。关于车驾，他们将四马战车献给太阳，双马战车献给月亮。但正如诗人所说：\Quote{厄里克托尼俄斯首先敢于将四马套上战车，并乘其车轮以胜利之速飞驰。}\OriginalTerm{厄里克托尼俄斯}{Erichthonius}，\OriginalTerm{伏尔甘}{Vulcan}和\OriginalTerm{弥涅耳瓦}{Minerva}之子，是地上不配激情之果实，是一个魔鬼怪物，不，是魔鬼本身，而不仅仅是一条蛇。但如果\Person{阿尔戈斯的特罗库斯}是首辆战车的制造者，他将那作品献给了\OriginalTerm{朱诺}{Juno}。如果\Person{罗慕路斯}在罗马首先展示四马战车，我想，他也在偶像中获得了一席之地，至少如果他等同于\OriginalTerm{奎里努斯}{Quirinus}的话。但既然战车有如此多发明者，车夫们自然也被染上了偶像崇拜的色彩；起初只有两种颜色，即白色和红色——前者献给晶莹白雪的冬季，后者献给赤红太阳的夏季；但后来，随着奢侈和迷信的发展，红色被一些人献给\OriginalTerm{玛尔斯}{Mars}，白色被另一些人献给\OriginalTerm{西风神}{Zephyrs}，而绿色献给大地母亲或春天，天蓝色献给天空和海洋或秋天。但既然每种偶像崇拜都被天主定罪，那献给自然元素的形式也必分担其谴责。\par

% structure: p0020 source=h2 target=section reason=relative-heading-rank; base=h2

\section{第十章}

现在让我们继续讨论戏院表演，我们已指出其与马戏团有共同起源，并带有同样的偶像崇拜名称——甚至从一开始它们就称为\Quote{Ludi}，且同样服务于偶像。它们在其盛况上也彼此相似：从圣殿和祭台到展示场所的同一游行，那可悲的乳香和鲜血的倾泻，伴以笛子和喇叭的音乐，全在占卜者和殡葬者——这两位负责丧葬仪式和祭献的可憎主管——的指挥下。因此，如同我们从\Quote{Ludi}的起源言及马戏团比赛，现在我们由此转向戏院表演，从展示场所开始。起初，戏院本身就是维纳斯的一座圣殿；简言之，正因如此，舞台表演才得以逃脱谴责并在世上立足。因为监察官们常常为了道德而首先压制新兴戏院，他们预见到戏院可能导致普遍的放荡；因此，从他们自身与我们的一致中，对外邦人已有一种见证，而在人类知识的预判中甚至是对我们观点的确证。故此，庞培大帝——仅次于他的戏院——在建立了那座万恶的堡垒后，担心有朝一日他的身后名受到监察官的谴责，便在其上方加建了一座维纳斯的圣殿；他通过公开宣告召集民众参与祝圣，并称之为不是戏院而是圣殿，\Quote{在此圣殿下，}他说，\Quote{我们安置了观看表演的阶梯座位。}于是他为一座常常受到谴责的建筑蒙上了面纱，并通过假装它为圣地，为一向该受责备的建筑遮掩；而借着迷信，他蒙蔽了良善纪律的眼睛。但维纳斯和巴克斯是亲密的盟友。这两个邪灵彼此立誓结盟，是醉酒和情欲的庇护者。因此维纳斯的戏院也是巴克斯的殿宇：因为他们恰当地也将\OriginalTerm{自由节}{Liberalia}的名称赋予其他戏院娱乐——这些娱乐除了被祝圣给巴克斯（正如希腊人的\OriginalTerm{狄俄尼索斯节}{Dionysia}），也是由他设立的；无疑，戏院的表演受这两位神祇的共同庇护。那种特别而独特地标志着舞台的放荡姿态和服饰是奉献给她们的——一位神祇因性别而淫荡，另一位因服饰而淫荡；而其声音、歌唱、琴瑟和笛子的服务属于阿波罗们、缪斯们、密涅瓦们和墨丘利们。基督徒啊，你必憎恨那些其作者必须是你完全憎恶对象的东西。因此我们现在想就戏院的技艺发表评论，也评论那些其作者的名字为我们所憎恶的东西。我们知道亡者的名字是虚无，他们的像也是虚无；但我们也十分清楚，当像被设立时，在这些名字下进行邪恶工作、对向它们表达的敬意洋洋得意、并假装神圣的，不是别的，正是被诅咒的灵，就是魔鬼。因此我们看到，技艺也被奉献给居住在其创始人名字中的存有；而且，那些发明者因其发现而在诸神中得有一席之地的东西，不能被视为没有偶像崇拜的玷污。不，就技艺而言，我们本应追溯到更早，并以这样的立场阻止一切进一步的争论：恶魔们从一开始就出于自身利益，在偶像崇拜的其他祸害中预先决定了公共表演的污秽，其目的是把人从他的主那里引开，将他束缚于自己的服务；他们通过赋予人表演所需的艺术天赋来执行其目的。因为除了他们自己，没有人会为他们心中的目标做准备和预备；他们也不会通过任何其他人将技艺赐予世界，而只通过那些在其名字、形象和历史中——他们为自身目标而设立奉献的把戏——的人。\par

% structure: p0022 source=h2 target=section reason=relative-heading-rank; base=h2

\section{第十一章}

为完成我们的计划，现在让我们继续考虑格斗比赛。它们的起源类似于游戏（\Emph{ludi}）。因此它们被奉为神圣或葬礼性质，因它们是为纪念外邦的偶像之神或亡者而设立。所以，它们被称为\OriginalTerm{奥林匹亚}{Olympian}以纪念朱庇特——在罗马称为卡比托利欧；\OriginalTerm{尼米亚}{Nemean}以纪念赫剌克勒斯；\OriginalTerm{地峡}{Isthmian}以纪念尼普顿；其余的称为\Emph{mortuarii}，属于死者。那么，如果偶像崇拜以亵渎的冠冕、祭司首领、属于各学院的随从，最后以其祭品的血来玷污战斗游行，有什么奇怪呢？关于\Quote{场所}再补充一句：在献给缪斯、阿波罗、密涅瓦以及献给马尔斯的技艺学院的共同场所中，他们以竞赛和号角声在竞斗场上仿效马戏团，而竞斗场是一座真正的圣殿——我指的是那位其节日得以庆祝的神祇的圣殿。体操技艺也起源于他们的卡斯托尔、赫剌克勒斯和墨丘利。\par

% structure: p0024 source=h2 target=section reason=relative-heading-rank; base=h2

\section{第十二章}

我们还需考察最有名、最受欢迎的\Quote{表演}。它被称为一种敬礼服务（\Emph{munus}），因它是一项职分，因为它的名称既是\Quote{\Emph{officium}}，也是\Quote{\Emph{munus}}。古人认为，在这庄严的礼仪中，他们向死者履行了职分；后来，以更精致的残忍，他们稍加改变了其性质。因为从前，他们相信亡者的灵魂可由人血平息，于是惯常购买俘虏或邪恶的奴隶，在他们的葬礼中献祭。后来，他们认为最好以欢喜的帷幕掩盖其不义。因此，凡他们为争斗而提供的人，然后尽力训练他们使用武器，只为让他们学会死亡，便在葬礼那天在墓地将其杀害。他们以凶杀来减轻死亡。这就是\Quote{Munus}的起源。但逐渐地，他们的精巧赶上了他们的残忍；因为这些人间猛兽找不到足够的精致快乐，除了看人被野兽撕碎的表演。那时安抚死者的供品被认为是属于葬礼祭献一类；而这些就是偶像崇拜：因为偶像崇拜事实上是对亡者的一种敬意；两者都是对死者的服务。再者，魔鬼居住在亡者的偶像中。再说到名称的问题，尽管这类表演已从对死者的尊荣转为对活人的尊荣——我是说，转为财务官和长官的职位，以及各种司祭职务——然而，既然偶像崇拜仍依附于尊荣之名，凡以其名义所做的，都沾染了其不洁。同样的评论也适用于\Quote{Munus}的游行，当我们看到它与这些尊荣本身相连的排场时；因为紫袍、权标、头带、冠冕，以及宣告和谕令，前日的圣宴，都离不开魔鬼的排场和魔鬼的邀约。那么，何必细述那恐怖之地，连伪证者的舌头也难以承受？因为竞技场奉献给比\Place{卡皮托利山}本身更多更可怕的名号，\Place{卡皮托利山}本是众魔鬼的庙宇。那里有多少人，就有多少不洁的邪灵。最后，关于在其中占有一席之地的艺术，我们只需说一句：我们知道它的两种娱乐是以\Person{玛尔斯}和\Person{狄安娜}为护佑者。\par

% structure: p0026 source=h2 target=section reason=relative-heading-rank; base=h2

\section{第十三章}

我想，我们已忠实地实现了我们的计划，展示了偶像崇拜之罪以多少不同的方式依附于表演，论及其起源、名称、装备、庆祝地点、技艺；我们可以毫无疑问地认为，对于我们这些已两次弃绝所有偶像的人来说，它们是全然不合适的。\BibleQuote{偶像算不得什么}（见：\BibleRef{格前 8:4}），正如宗徒所说，而是他们献上的敬意是献给魔鬼，魔鬼才是这些被祝圣的偶像的真正居住者，无论这些偶像是死人还是（他们以为的）神明。因此，因为他们有共同的根源——他们的死者和他们的神祇是一体的——我们戒绝这两种偶像崇拜。我们对庙宇的厌恶不亚于对坟墓：我们与两边的祭台都无关，我们不敬拜任何偶像；我们不向神明献祭，也不为亡者献葬礼供物；不仅如此，我们也不吃这两方所献的物，因为我们不能既参加天主的筵席又参加魔鬼的筵席。（见：\BibleRef{格前 10:21}）如果我们使喉咙和肚腹远离这些玷污，那么我们岂不更应当持守我们更高贵的部分，即我们的耳朵和眼睛，远离那些偶像崇拜和葬礼的娱乐？这些娱乐不经过身体，却直接在精神与灵魂中消化，而精神与灵魂的纯洁，天主有权从我们这里要求，远胜于我们身体器官的纯洁。\par

% structure: p0028 source=h2 target=section reason=relative-heading-rank; base=h2

\section{第十四章}

既然我们已经充分确立了偶像崇拜的指控，单凭这一点就足以成为我们放弃表演的理由，让我们现在\Emph{ex abundanti}以另一种方式来看这个主题，尤其是为了那些安慰自己认为我们所敦促的节戒并非明文规定的人，仿佛在谴责世俗私欲的同时，并没有包含反对所有这些娱乐的充分宣告。因为正如有对金钱、地位、饮食、不洁享乐或荣耀的贪欲，也有对快乐的贪欲。而表演正是一种快乐。因此我认为，在贪欲的总称下，快乐被包含在内；同样，在快乐的总概念下，你们有特定的一类，就是\Quote{表演}。但我们已说过展览场所的情况：它们本身并不污染，而是由于在其中所行的事——它们从中吸收不洁，然后又喷溅到别人身上。\par

% structure: p0030 source=h2 target=section reason=relative-heading-rank; base=h2

\section{第十五章}

既已如所述，足够处理了那主要论点——即一切表演都染有偶像崇拜的污秽——并充分论述了这一点，如今让我们将表演的其他特性与天主的事物相对照。天主命我们以平静、温和、安静、平安与圣神相处，因为这些事唯独符合祂本性之善、祂的温柔和敏感，且不以愤怒、恶念、怒气或忧伤使祂忧烦。那么，这如何能与表演相协调呢？因为表演总是引致精神的激动，因为哪里有快乐，哪里就有强烈的情感给快乐增添滋味；哪里有强烈的情感，哪里就有竞争，轮流给情感增添滋味。再者，哪里有竞争，哪里就有愤怒、苦毒、怒气与忧伤，以及从它们流出的一切恶事——这些完全与基督的宗教不合。因为即使有人以合乎其身份、年龄或本性的适度方式享受表演，他内心仍非毫无扰动，没有一些未言明的内在骚动。没有人能毫无强烈激动地参与这样的快乐；也没有人能在其激动之下而不有自然的失足。这些失足又产生贪欲。若没有贪欲，就没有快乐；而若无所获却仍前往，此人便应受轻浮之责；依我看，连这也与我们无关。此外，一个人仅仅置身于那些他因不愿相似而承认与之无同情心的人中间，便宣告了自己的定罪。我们自己不行这样的事还不够，除非也与行这样事的人断绝一切关联。圣经说：\Quote{你若见盗贼，便与他同谋。}但愿我们甚至不与这些恶人共居同一个世界！但这愿望虽无法实现，然而如今我们在属世之事上已与他们分离；因为世界是天主的，但属世的是魔鬼的。\par

% structure: p0032 source=h2 target=section reason=relative-heading-rank; base=h2

\section{第16章}

既然一切的激情冲动都被禁止，我们便被隔绝于各种景观，尤其是竞技场，在那里这样的激动如同在其固有元素中般主导。看那人群来到时已情绪激昂、已喧闹、已盲目于激情、已因赌注而骚动。裁判官对他们来说太慢了；他们的眼睛不停地转动，仿佛与瓮中的签一同转动；然后他们都急切地悬望着信号；那里有一阵共同疯狂的齐声呼喊。看他们如何因愚昧的言语而\Quote{魂不守舍}。他们喊着：\Quote{他扔了！}并各自向邻人宣布众人所见之事。我有最清楚的证据表明他们的盲目；他们看不见真正被扔下的是什么。他们以为那是一块\Quote{信号布}，但那却是魔鬼从高处被掷下的形象。结果便是，他们陷入愤怒、激情、纷争，以及一切献身于和平的人绝不应放纵的事。随后有咒骂和指责，却无仇恨的原因；有喝彩声，却无值得赞美之处。参与这一切的人——不能自制——为自己获得了什么？除非也许获得了那使他们不能自制的东西：他们因别人的悲哀而悲伤，因别人的喜乐而欢喜。他们所欲或所憎的一切，完全与自身无关。因此他们的爱是无益的，恨是不义的。或者，无缘无故的爱是否比无缘无故的恨更为合法？天主确实禁止我们甚至在有理由恨时去恨，因为祂命令我们爱仇敌。天主禁止我们咒骂，即使有些原因，祂命令我们祝福那咒骂我们的人。还有什么比竞技场更冷酷无情呢？在那里人们甚至不宽恕自己的统治者和同胞公民。如果其中的任何疯狂在其他地方对天主的圣人是合宜的，那么在竞技场中也是合宜的；但如果它们在任何地方都不对，那么在那里也不对。\par

% structure: p0034 source=h2 target=section reason=relative-heading-rank; base=h2

\section{第17章}

同样，我们不是奉命要摒除一切不贞洁吗？正因如此，我们又不得进入剧场——那不贞洁特有的居所，那里除了在别处可耻之事外，别无他物受人推崇。所以，欲得其神的最大恩宠，最佳途径便是卑劣：阿特拉闹剧所表演的，扮女装的小丑所展示的，毁灭一切天然的廉耻，以致人们在家比在戏中更容易脸红，最终，从童年起便在默剧演员身上施行，以使他成为伶人。就连那些妓女——公共情欲的牺牲品——也被带上舞台，加重她们的悲惨，因为在那里，她们当着同性的面，而她们平日只向同性遮掩自己；她们被公开示众于各种年龄与各种等级面前——她们的居所、她们的收益、她们的赞誉，都被展示出来，甚至是在那些不该听闻这些事的人耳中。其他事情，我暂不提及，那些事最好隐没于它们自身的黑暗和幽深的洞穴中，以免玷污白昼之光。让元老院，让所有等级，因极度羞耻而脸红！为何连这些可怜的女子，她们用自己的动作摧毁了自己的廉耻，畏惧白昼之光与民众的目光，至少每年一度还知羞耻。但如果我们应当憎恶所有不贞洁的事，那么，有什么理由可以听那些我们不可说的事呢？因为一切放肆的言语，甚至一切闲话，都被天主定罪。同样，为什么可以看那些做起来可耻的事呢？那些从口里出来污秽人的事，怎么不认为从眼耳进入时也污秽人呢？——眼耳是灵的贴身仆役——而仆役不洁的灵，绝不可能是纯洁的。所以，在不贞洁的禁止中，剧场也被禁止了。再者，如果我们轻视世俗文学，视之为天主眼中的愚拙，那么对于那些源自世俗文学的悲剧或喜剧表演，我们的职责就十分明显了。如果悲剧和喜剧是血腥与纵欲、不敬与放荡的罪恶与情欲的发明者，那么，连回忆那些残暴或卑劣之事都是不好的。你在行为上拒绝的事，也不该在言语上欢迎。\par

% structure: p0036 source=h2 target=section reason=relative-heading-rank; base=h2

\section{第十八章}

但如果你争辩说，圣经中提到了赛马场，我立刻承认。但你不会拒绝承认，那里所做的事是你不可观看的：打击、踢踹、耳光，以及一切手的鲁莽，以及任何像这样对人面容的毁损——那无异于对天主肖像的毁损。你绝不会认可那些愚蠢的赛跑和投掷技艺，以及更愚蠢的跳跃；你不会从伤害性或无用的力量展示中找到乐趣；你肯定不会认可那些追求人为身体、企图超越造物主之工的努力；你对那些在和平的闲散中滋养出来的希腊运动员，只会抱有完全相反的态度。而摔跤的技术是魔鬼的东西。魔鬼与最初的人类摔跤，并将他们摔死。摔跤的姿势本身就带有蛇类的力量：紧缠不放——折磨般地搂抱——滑溜地溜走。你不需要冠冕；为何你要从冠冕中寻求快乐呢？\par

% structure: p0038 source=h2 target=section reason=relative-heading-rank; base=h2

\section{第十九章}

现在我们来看圣经如何谴责竞技场。如果我们能认为沉溺于残忍、不敬与凶暴是正当的，那我们就去吧。如果我们如人们所说的那样，就让我们在那里用人血来宴乐。当然，让有罪者受罚是好的。除了罪犯本人，谁会否认呢？然而，无辜者不会从他人的苦难中找到快乐：他反而悲伤，因为一个兄弟犯了如此重罪，以致需要如此可怕的惩罚。但谁能向我保证，总是有罪者被判决喂野兽或受其他刑罚，而无辜者从不会因法官的报复、辩护的软弱或刑讯的压力而受苦呢？那么，我最好还是对恶人所受的惩罚一无所知，免得我不得不也知道好人遭遇不测——如果这也能算作好人的话！至少，没有被控罪名的角斗士被卖来供表演之用，成为公众娱乐的牺牲品。即使是那些被司法判决送往竞技场的人，多么怪异啊，他们在受刑时，只因较轻微的过失，就沦为杀人犯！但我这些话是对异教徒说的。至于基督徒，我不再额外多言，以免侮辱他们，说明他们应当对此类表演持何等厌恶的态度；然而，没有人比我更能完整地阐述整个主题，除非是那些仍然习惯于观看表演的人。我宁愿不完满，也不愿启动回忆。\par

% structure: p0040 source=h2 target=section reason=relative-heading-rank; base=h2

\section{第二十章}

那么，那些为了不愿失去享乐而声称我们无法指出禁止这项节制的具体字句或确切地方的人，其推理是多么虚妄——不，是何等绝望——啊！他们还说天主的仆人被直接禁止与这类集会有关联。我最近听到一位爱看戏的人为自己所作的奇特辩护。他说：\Quote{太阳，不，是天主自己，从天上看这演出，并不染污秽。}的确，太阳也将光线倾注进公共下水道而不受玷污。至于天主，但愿一切罪过都瞒过祂的眼，好让我们都能逃脱审判！但祂也看见抢劫；祂看见谎言、奸淫、欺诈、偶像崇拜，以及这些演出；正因如此，我们才不看它们，免得那全视者看见我们。人啊，你把罪犯和审判者等同起来了：罪犯因被看见而成为罪犯，审判者因看见而成为审判者。难道我们注定要在马戏场围墙外发疯吗？在剧场门外我们执意淫乱，在跑道外执意傲慢，在竞技场外执意残忍，因为门廊、看台和幕布之外，天主也有眼目吗？凡天主永远谴责的事，在任何地方、任何时候都不能幸免于罪责；在任何地方、任何时候都不可以做那些你无论在何时何地都不允许做的事。真理之自由在于它不变的意见和判断，这构成了它的圆满，使它有权获得完全的统治、不变的敬畏和忠信的服从。真正善或真正恶的事物不能是别的。但在一切事上，天主的真理是不变的。\par

% structure: p0042 source=h2 target=section reason=relative-heading-rank; base=h2

\section{第二十一章}

外邦人没有真理的完整启示，因为他们不受天主的教导，他们随从自意和情欲来判定善恶，使在此处为善的在彼处为恶，在彼处为恶的在此处为善。于是奇怪地发生这样的事：同一个人，在公开场合连出于自然需要而撩起外衣都几乎做不到，却在马戏场中脱下它，仿佛要当众暴露自己；那位小心翼翼保护并守护自己女儿贞洁的耳朵远离一切污言秽语的父亲，亲自带她去剧场，让她暴露在一切卑鄙的言语和姿态之中；那个在街头上对争吵的拳击手动手或辱骂的人，在竞技场中却对更严重的打斗给予一切鼓励；那个对死于自然规律之下的尸体感到恐惧的人，在竞技场中却以最大的耐心凝视着那些被撕裂、被扯碎、沾染着自己鲜血的躯体；不，正是那个来看演出的人，因为他认为杀人犯应该为罪行受罚，却用棍棒和鞭子驱使不情愿的剑斗士去杀人；那个要求用狮子来处决每一名更凶残的杀人犯的人，却为野蛮的剑士提供杖，并用自由之帽奖励他。是的，他还要求把可怜的牺牲者送回来，好让他看一眼他的脸——饶有兴趣地就近审视那个他本想从远处撕成碎片的人：如果他没有这个愿望，那他就更加残忍了。\par

% structure: p0044 source=h2 target=section reason=relative-heading-rank; base=h2

\section{第二十二章}

这有什么奇怪的呢？这种矛盾正是我们从那些因感觉无常和判断善变而混淆善恶本性的人身上所预料到的。为什么，演出的创办者和管理者们在赞颂车夫、演员、摔跤手和那些最可爱的剑斗士（人们为他们出卖灵魂，女人甚至出卖肉体）的同时，却轻视和践踏他们，尽管为了他们，人们也犯了所谴责的行为；不，他们甚至判定他们不名誉并剥夺其公民权利，将他们排除在\OriginalTerm{库利亚}{Curia}、\OriginalTerm{演讲台}{rostra}、元老院和骑士阶层以及所有其他荣誉和某些特权之外。这是何等颠倒！他们从那些他们惩罚的人身上得到快乐；他们贬低那些他们同时给予认可的人；他们赞美技艺却谴责艺术家。以那些使他们眼中的人有价值的事情来污蔑他们，这是何等荒谬！不，这是一份何等明显的供认，即这些事情是恶的，因为它们的作者即使处于最高的恩宠之中，也难免带着耻辱的标记！\par

% structure: p0046 source=h2 target=section reason=relative-heading-rank; base=h2

\section{第二十三章}

于是，人自身的反省，即便在享乐之甘甜中，仍促使他思考：那些如此行事的人应被判定遭受不名誉的厄运，丧失与拥有生活尊位相关的一切利益。那么，天主的正义岂不更惩罚那些投身于这些技艺的人吗？天主会喜悦那个驭车者吗？他搅扰如此多的灵魂，激起如此多的狂怒之情，制造如此多样的心态——或如司祭般戴冠，或如皮条客般着色，被魔鬼装扮，好让他在战车上被卷走，仿佛要带走厄里亚一般？天主会喜悦那个对自己动剃刀、完全改变面容的人吗？他不尊重自己的脸，不满足于将其尽可能变得像农神、伊西斯和巴克斯，却静静地将它交给侮辱性的打击，仿佛在嘲弄我们的主？的确，魔鬼也将其纳入他的教训：脸颊要温顺地交给掌掴者。同样，他用高底鞋使悲剧演员增高，因为\BibleQuote{没有人能增加自己的身量一肘。}\BibleRef{玛 6:27} 他的愿望是使基督成为说谎者。至于戴面具，我问：这是符合禁止制造任何肖像、尤其是禁止制造按自己肖像所造之人的天主的心意吗？真理的创造者憎恶一切虚假；祂视一切不真实为奸淫。因此，祂谴责各种形式的虚伪，绝不会认可任何扮演声音、性别或年龄的行为；绝不会认可伪装的爱情、愤怒、呻吟和眼泪。此外，正如在祂的法律中宣告，穿戴女性服装的人是可咒骂的，\EditorialNote{申命纪第廿二章} 那么，祂对那个甚至被培养成扮演女子的哑剧演员会作何判决？拳击手难道能逃脱惩罚吗？我想，他是在受造时就得到了这些拳击带伤痕、拳头的厚皮和耳朵上的增生物吧？天主赐给他眼睛，无非是为了让它们在战斗中被打掉！我且不提那个为了救自己而把别人推进狮子之路的人，免得他在竞技场上杀死同一个人时，杀人罪过不够重。\par

% structure: p0048 source=h2 target=section reason=relative-heading-rank; base=h2

\section{第二十四章}

我们还要用多少其他方式进一步证明，与竞技表演相关的一切，没有一件是天主所赞许的，或不经祂赞许而适合天主的仆人？如果我们已经成功阐明，它们完全是为魔鬼的缘故而设立，并且完全是用魔鬼的事物来筹办（因为凡不是天主的，或不蒙祂悦纳的，都属祂邪恶的对手），那么这仅仅意味着，在这些表演中，你拥有那魔鬼的盛装，即我们在信德\OriginalTerm{印记}{seal}中所弃绝的。我们不应与我们所弃绝的事物有任何关联，不论是在行为或言语上，不论是通过观看它们还是期待它们；然而，当我们不再为那圣洗的誓愿作证时，难道我们不是弃绝并撤销了它吗？那么，我们是否还需向外教人请教？让他们告诉我们，基督徒经常出入表演是否正当。为什么，拒绝这些娱乐是他们判断一个人接受了基督信仰的首要标志。那么，若有人抛弃了信德的独特标记，他显然就是否认了它。一个这样做的人，你还能对他保留什么希望？当你投奔敌人的营地时，你放下武器，遗弃旗帜和对统帅的效忠誓言：你的生死都与新同伴拴在一起。\par

% structure: p0050 source=h2 target=section reason=relative-heading-rank; base=h2

\section{第二十五章}

坐在没有天主的地方，人还会思念他的造物主吗？当灵魂为了一个驭车者而激烈争吵时，还能有平安吗？在狂热的兴奋中被激起，他会学会谦逊吗？不，在整个事件中，他遇到的最大诱惑莫过于男女的艳丽装束。情感的相互交织、在偏爱问题上彼此的一致与分歧——在如此亲密的共处中——都会点燃情欲的火花。此外，去观看表演几乎别无其他目的，除了看与被看。当悲剧演员在朗诵时，人还会想到先知性的劝诫吗？在柔媚演员的节拍中，他还能唤起一首圣咏吗？当运动员激烈搏斗时，他还能宣称不可还击吗？当目光盯着熊的撕咬和网斗士的网时，他还能被怜悯之心打动吗？愿天主从祂的子民中除去任何对残忍享受的这种狂热渴望！因为从天主的教会到魔鬼的场所——从天穹到猪圈，如俗语所说——是多么可怕：向天主举起双手，然后又因对演员的掌声而疲倦；从口中发出对圣体宣念的\Quote{阿们}，又为角斗士作见证；向任何人喊\Quote{永远}，而不是向天主和基督！\par

% structure: p0052 source=h2 target=section reason=relative-heading-rank; base=h2

\section{第二十六章}

那些进入表演诱惑的人，为何不能也接近邪灵？我们有一个妇女的例子——主自己作证——她去剧院，回来时被附了魔。所以在驱逐时，当那不洁之物被斥责竟敢攻击一个信徒时，它坚定地回答：\Quote{的确，我最公义地做了这事，因为我在我的领地找到了她。}另一个著名的案例：一个妇女听过一个悲剧演员的演唱，当夜她在梦中看见一块麻布——演员的名字同时被提及并受到强烈谴责——五天后，那妇女就死了。又有多少其他确凿的证据，是我们从那些因在表演中与魔鬼为伴而离弃了主的人身上得到的！因为\BibleQuote{没有人能服侍两个主人。}\BibleRef{玛 6:24} \BibleQuote{光明与黑暗有什么相通？生命与死亡有什么相交？}\BibleRef{格后 4:14}\par

% structure: p0054 source=h2 target=section reason=relative-heading-rank; base=h2

\section{第二十七章}

我们应当憎恶这些外教人的聚会和集会，即使没有别的理由，也因在那里天主的名字被亵渎——因那里每天向我们高喊\Quote{丢给狮子！}——因迫害的法令常从那里发出，诱惑被散播。如果你被卷入那股不敬裁判的汹涌洪流中，你将如何？并非那里有人会伤害你：没有人知道你是基督徒；但思想一下你在天堂的遭遇。因为正当魔鬼在教会中肆虐时，你难道怀疑天使正从天上下望，注视每一个人，谁说话，谁听亵渎之言，谁献出舌头，谁献出耳朵为撒殚服务来对抗天主吗？难道你不该避开那些基督敌人聚集的阶梯，那个一切瘟疫的座位，以及那被邪恶呼声玷污的笼罩其上方的空气吗？就算你在那里有愉快的事，本身令人愉悦且无罪的事，甚至有些优秀的事。没有人用苦胆和藜芦稀释毒药：可咒之物被放入调味精美、味道最甜的佐料中。同样，魔鬼在他调制的致命毒液中，也放入最令人愉悦、最可接纳的天主的事物。因此，那里的一切，无论是勇敢、高贵、响亮、悦耳或滋味精美，都只当作毒饼上的蜜滴；不要过于看重你对它的喜好的味道，而忽略你因它的诱惑所冒的危险。\par

% structure: p0056 source=h2 target=section reason=relative-heading-rank; base=h2

\section{第二十八章}

让魔鬼的宾客用这些美馔饱享。地方、时间，以及邀请者，都是他们的。我们的宴席，我们的婚宴，尚未来临。我们不能与他们一同坐席，正如他们也不能与我们一同坐席。此事轮流转。如今他们欢乐，而我们忧愁。耶稣说：\Quote{世界要欢乐，你们却要忧愁。}\BibleRef{若 16:20} 那么，我们当在外教人欢乐时哀悼，好使我们在他们忧愁的日子欢乐；免得现在分享他们的欢乐，那时也分享他们的悲伤。基督徒啊，你太娇气了，竟想今生和来世都享快乐；不，你真是愚昧，若以为今生的快乐是真正快乐。例如，哲学家称安静和安息为快乐；他们以此为福乐，以此为消遣；甚至以此夸耀。你竟渴望终点、舞台、尘土和竞技场！我要你回答我：我们岂不能没有快乐而活，却不能不快乐而死吗？因为我们的愿望岂非宗徒的愿望，离开世界，被接去与主同在吗？\BibleRef{斐 1:23} 你的快乐在哪里，你的渴慕就在哪里。\par

% structure: p0058 source=h2 target=section reason=relative-heading-rank; base=h2

\section{第二十九章}

即使如今，你若想在今生度日于享乐中，你何以如此忘恩，竟以为天主所赐予你的许多精致快乐不够？难道要你感恩地承认吗？因为还有什么比天主圣父和我们的主与我们相和，比真理的启示，比承认我们的过错，比赦免我们过去生命中无数的罪，更令人喜悦？有什么比厌恶快乐本身、轻视世界所能给的一切、真正的自由、纯洁的良心、知足的生活、以及免于对死亡的一切恐惧，更伟大的快乐？有什么比践踏外邦人的神祇、驱逐邪灵、施行治愈、寻求神圣的启示、为天主而活，更高贵？这些才是快乐，这些才是相称于基督徒的景观——神圣、永恒、自由。把它们算作你的赛马场，定睛看世界的赛程、流逝的季节，计算时间的时期，渴望最终圆满的终点，捍卫教会的团体，为天主的号令而惊起，为天使的号角而振奋，在殉道的棕榈枝上夸耀。如果剧院的文学使你喜悦，我们自己有丰富的文学——大量的诗歌、格言、歌曲、箴言；这些不是神话，而是真理；不是艺术的花招，而是朴素的现实。你也想要打斗和摔跤吗？这些也不缺乏，而且不是小事。看贞洁战胜不洁，忠信杀死背信，怜悯击打残忍，谦虚使无耻黯然失色：这些是我们中间的竞赛，在这当中\Emph{我们}赢得冠冕。你也想要一点血吗？你有基督的血。\par

% structure: p0060 source=h2 target=section reason=relative-heading-rank; base=h2

\section{第三十章}

然而，我们的主那即将来临的显现，是何等的景象！如今被万民承认，如今备受尊崇，如今是凯旋者！天使军旅的欢腾是何等景象！复活的圣人们的荣耀是何等景象！此后义人的国度是何等景象！新耶路撒冷城是何等景象！是的，还有别的景象：那最后审判的日子，带着它永恒的结局；那被万民所不预料、成为他们讥嘲主题的日子，当世界因年迈而苍白，连同它的一切产物，都将在同一场大火中烧尽！那时，何等浩大的景象涌入眼帘！那里有什么激起我的钦佩？有什么引起我的讥嘲？哪一幕景象给我喜乐？哪一幕激励我欢腾？——当我看见如此众多显赫的君王，他们被接升天曾公开宣告，如今却与伟大的\Person{朱庇特}本人在最深的黑暗中呻吟，还有那些曾见证他们欢腾的人；此外，各省的总督，他们迫害基督徒，在比他们昔日骄傲中对基督信徒所发之火更猛烈的火焰中。还有世上的智者，事实上那些哲学家，他们教导门徒说天主毫不关心属下的任何事物，并惯于向他们保证：要么他们没有灵魂，要么灵魂永远不会回到死亡时离开的身体，如今在那些被欺骗的可怜人面前羞愧难当，同一火焰吞噬了他们！诗人们也在颤抖，不是在\Person{拉达曼托斯}或\Person{米诺斯}的审判座前，而是在意想不到的基督面前！那时我将有更好的机会聆听悲剧演员，他们在自身的灾难中声音更响；观看戏剧演员，在消融的火焰中更加\Quote{放荡}；察看战车驭手，在他火战车中通体炽热；目睹角斗士，不在他们的体育馆中，而在火浪中翻滚；除非到那时我仍不屑于关注这些罪的仆役，而是热切渴望将那永不满足的凝视锁定在那些向主发泄怒气的人身上。我将说：\Quote{这个，}我将说，\Quote{这就是那个木匠或雇工的儿子，那个安息日破坏者，那个撒玛黎雅人和附魔者！这就是你们从\Person{犹达斯}那里买来的那位！这就是你们用苇秆和拳头击打、轻蔑地吐唾沫、给他喝苦胆和醋的那位！这就是他的门徒偷偷偷走、好让人说他已复活了的那位，或者是园丁搬走了，免得他的莴苣被来访的人群损坏的那位！}有哪位财务官或司铎，凭其慷慨，会将看见并欢庆如此之事的恩惠赐予你？然而，即使现在，我们也凭信德在想象的描绘中多少拥有它们。但是，那些目所未见、耳所未闻、人心未曾想到的事情，到底是什么呢？无论它们是什么，我相信它们比竞技场、两个剧场和所有赛马场更高贵。\par

\SourceRecord{Translated by S. Thelwall.；Ante-Nicene Fathers；Vol. 3.；Edited by Alexander Roberts, James Donaldson, and A. Cleveland Coxe.；Buffalo, NY: Christian Literature Publishing Co.,；1885.；<http://www.newadvent.org/fathers/0303.htm>.}
